\section{Rounds of Adaptivity}\label{app:D}
\subsection{Rounds of Adaptivity}\label{app:roundshier}
\paragraph{Interviews with $r$ rounds.}
An $r$-round interview chooses coordinate batches
$B_1,\;B_2(x_{B_1}),\;\ldots,\;B_r(x_{B_1\cup\cdots\cup B_{r-1}})$.
Each batch is a function of the answers to earlier batches only, and
$\sum_s|B_s|\le b$ on every input. Let $\A_r(b)$ be the finite menu of
deterministic $r$-round interviews and $V_r(d,b)=V(d,\A_r(b))$; randomized
interviews are mixtures and are already represented by the design
probabilities in \eqref{eq:unbiasedprogram}. Fixed forms are $\A_1(b)$.
Sequentializing each batch shows that every $r$-round interview is a depth-$b$
tree, and a depth-$b$ tree is a $b$-round interview with one query per round.
Hence $\A_1(b)\subseteq\A_2(b)\subseteq\cdots\subseteq\A_b(b)=\A_{\rm branch}(b)$
and $V_1\ge V_2\ge\cdots\ge V_b=V_{\rm branch}$. All menus identify the same
contrasts, $\deg(d)\le b$, because the two extreme menus do
(Appendix~\ref{app:id}). Theorems~\ref{thm:finite} and~\ref{thm:asym} apply to
each finite menu $\A_r(b)$, so the lower bounds below hold against biased,
sequential and label-dependent procedures.

\begin{theorem}[Rounds of adaptivity, full statement of Theorem~\ref{thm:rounds}]\label{thm:roundsfull}
\textup{(a)} For $1\le r\le b\le m$ and every nonzero contrast of degree at most $b$ on $\{-1,1\}^m$,
\[
 V_{\rm branch}(d,b)\le V_r(d,b)\le
 \min\Big\{2^{b-r},\binom{m-r}{b-r}\Big\}\,V_{\rm branch}(d,b).
\]
\textup{(b)} For integers $r,\kappa\ge1$ put $D=2^\kappa$ and $b=r\kappa+1$.
The pointer-chasing contrast $N_{r,\kappa}$ defined below satisfies
\[
 V_{r+1}(N_{r,\kappa},b)=V_{\rm branch}(N_{r,\kappa},b)=1,\qquad
 \frac{D}{b}\le V_r(N_{r,\kappa},b)\le D.
\]
\end{theorem}

\begin{corollary}[Two rounds can be exponentially worse than branching]\label{cor:tworounds}
For every $b\ge3$,
\[
 \sup_{m,d}\frac{V_2(d,b)}{V_{\rm branch}(d,b)}\ \ge\ \frac{2^{\lfloor(b-1)/2\rfloor}}{b},
 \qquad
 \sup_{m,d}\frac{V_r(d,b)}{V_{r+1}(d,b)}\ \ge\ \frac{2^{\lfloor(b-1)/r\rfloor}}{b}
 \quad(1\le r<b).
\]
\end{corollary}
\begin{proof}
Take $\kappa=\lfloor(b-1)/r\rfloor\ge1$ and $N_{r,\kappa}$, which uses
$r\kappa+1\le b$ queries along every path. Branching and $(r+1)$-round
interviews still compute it exactly under the larger cap $b$, while the hit
bound in the proof of part (b) uses only the total number $b$ of queries.
For the first display use $r=2$ and $V_{\rm branch}\le V_3$.
\end{proof}

The constant separation $V_2\ge4/3$ of Corollary~\ref{cor:roundseparation}
is therefore not the end of the hierarchy. Part (b) and the corollary place
the exponent for $2\le r<b$ between $\lfloor(b-1)/r\rfloor-\log_2b$ and $b-r$;
the binary pointer tree (Theorem~\ref{thm:binarytree}) raises the lower end to
$b-1-r\log_2(1+2b/r)$, its products do so in polynomial dimension, and
Appendix~\ref{app:pointerexact} computes the constants of the tree at small caps.

\paragraph{Proof of part (a).}
The left inequality is menu inclusion. For the right inequality pad a
depth-$b$ tree $\tau$ as in Appendix~\ref{app:universal}, so that every path
has $b$ distinct coordinates. In rounds $1,\ldots,r-1$ ask one question per
round and follow $\tau$ exactly. This reaches a node $v$ of depth $r-1$,
whose padded subtree has depth $b-r+1$. Its path sets all contain the subtree
root and are determined by the first $b-r$ answers below $v$; hence their
number $L_v$ is at most $2^{b-r}$, and also at most $\binom{m-r}{b-r}$
because the remaining $b-r$ coordinates are chosen from the $m-r$
coordinates not on the followed prefix or at $v$. In round $r$ choose one of
these $L_v$ sets uniformly, independently of the person, and attempt to
complete the path. On completion return $L_v$ times the tree score and
otherwise zero. Conditionally on $x$ the mean equals the tree score and the
second moment is $L_v$ times its square. The mixture and merging argument of
Appendix~\ref{app:universal} applies verbatim to randomized branching
designs, so every feasible branching design is simulated with second
moment inflated by at most $\max_vL_v$.

\paragraph{The pointer-chasing contrast.}
Levels $\ell=1,\ldots,r$ consist of registers. Level $\ell$ has $D^{\ell-1}$
registers $R_\ell(i_1,\ldots,i_{\ell-1})$, indexed by
$(i_1,\ldots,i_{\ell-1})\in[D]^{\ell-1}$, and each register stores $\kappa$
binary coordinates encoding an element of $[D]$. There are $D^r$ data
coordinates $Z_{i_1\cdots i_r}$. The pointer path is $j_1=R_1$ and
$j_\ell=R_\ell(j_1,\ldots,j_{\ell-1})$, and
\[
 N_{r,\kappa}(x)=Z_{j_1\cdots j_r}
 =\sum_{i\in[D]^r}Z_i\prod_{\ell=1}^r\mathbf 1\{R_\ell(i_{<\ell})=i_\ell\}.
\]
Each register indicator has degree $\kappa$, so $\deg N_{r,\kappa}=r\kappa+1=b$.
Reading the $\kappa$ coordinates of the level-$\ell$ register on the pointer
path in round $\ell$, and the selected data coordinate in round $r+1$,
computes the contrast with $b$ queries in $r+1$ rounds. With Jensen's
inequality this gives $V_{r+1}=V_{\rm branch}=1$. For the upper bound
$V_r\le D$, read the pointer registers of levels $1,\ldots,r-1$ in the first
$r-1$ rounds. In round $r$, read the level-$r$ register together with one
data coordinate $Z_{j_1\cdots j_{r-1}J}$, where $J$ is uniform on $[D]$ and
drawn in advance. The score $DZ_{j_1\cdots j_{r-1}J}\mathbf 1\{J=j_r\}$ is
unbiased with second moment $D$ at every profile.

\paragraph{The lower bound $V_r(N_{r,\kappa},b)\ge D/b$.}
Let $\U$ be the uniform law and write $c^\star=(j_1,\ldots,j_r)$ for the
selected data cell. Since $\E_\U[N^2]=1$, the dual formula~\eqref{eq:vdual}
with witness $h=N$ gives $V_r\ge1/\max_{\iota\in\A_r(b)}E_\iota$, where
$E_\iota=\E_\U[\E_\U[N\mid\ell_\iota(X)]^2]$. Fix a deterministic $r$-round
interview $\iota$.

\emph{Step 1: a generous transcript.}
Let $\iota^+$ ask the same questions as $\iota$, but at the end of every
round additionally reveal the full content of every register touched in that
round. The transcript of $\iota$ is a function of that of $\iota^+$, and
conditional expectation onto a larger $\sigma$-field cannot decrease the
second moment, so $E_\iota\le E_{\iota^+}$.

\emph{Step 2: only the selected data coordinate carries signal.}
Let ${\rm Obs}$ be the set of data coordinates queried by $\iota^+$. For a
data cell $c$, on the transcript event $\{c\notin{\rm Obs}\}$ the transcript
and the pointer path are functions of the coordinates other than $Z_c$,
because the interview never reads $Z_c$. Since $Z_c$ is an independent
uniform sign, $\E[Z_c\mathbf 1\{c^\star=c\}\mid\ell_{\iota^+}]$ vanishes there.
Hence
\[
 \E[N\mid\ell_{\iota^+}]=\sum_{c\in{\rm Obs}}Z_c\,\PP(c^\star=c\mid\ell_{\iota^+}),
\]
so $|\E[N\mid\ell_{\iota^+}]|\le\PP(c^\star\in{\rm Obs}\mid\ell_{\iota^+})\le1$ and
\[
 E_{\iota^+}\le\E\big[\PP(c^\star\in{\rm Obs}\mid\ell_{\iota^+})\big]
 =\PP(c^\star\in{\rm Obs}).
\]

\emph{Step 3: observing $c^\star$ in $r$ rounds requires a lucky guess.}
Let $\G_t$ be the $\sigma$-field of the $\iota^+$ transcript before round $t$,
and let $K_t$ be the largest $k$ such that the pointer registers
$R_1,R_2(j_1),\ldots,R_k(j_{<k})$ are all revealed before round $t$; thus
$K_1=0$ and $K_t$ is nondecreasing. A query of round $t$ is \emph{deep} if it
touches a register of level at least $K_t+2$ or a data coordinate while
$K_t\le r-1$; it is a \emph{hit} if, in addition, the touched register or
data coordinate lies on the pointer path. Every deep location is indexed by a
tuple whose $(K_t+1)$st entry is $\G_t$-measurable, and lies on the pointer
path only if this entry equals $j_{K_t+1}$. The latter is the content of the
register $R_{K_t+1}(j_{\le K_t})$, which is unrevealed before round $t$.
For an adaptive sequence of queries to independent uniform coordinates, the
unrevealed coordinates remain independent and uniform given the
transcript. Hence $j_{K_t+1}$ is uniform on $[D]$ given $\G_t$, and every deep
query is a hit with conditional probability at most $1/D$.

If no hit occurs in rounds $1,\ldots,t$, then $K_{t+1}\le K_t+1$. Indeed,
$K_{t+1}\ge K_t+2$ requires the pointer register of level $K_t+2$ to be
revealed by the end of round $t$, say in round $t'\le t$. Since
$K_{t'}\le K_t$, that query was deep in round $t'$ and lay on the pointer path,
so it was a hit. By induction, without hits $K_t\le t-1\le r-1$ for every
$t\le r$. Then every query of a data coordinate in rounds $1,\ldots,r$ is deep,
and $c^\star$ can be queried only by a hit. Since the total number of queries
is at most $b$,
\[
 \PP(c^\star\in{\rm Obs})\le\PP(\text{some hit})
 \le\E\Big[\sum_{t=1}^r\sum_{\text{deep queries of round }t}\PP(\text{hit}\mid\G_t)\Big]
 \le\frac bD .
\]
Combining the three steps gives $E_\iota\le b/D$ for every $\iota\in\A_r(b)$,
and therefore $V_r\ge D/b$. \hfill$\square$

\paragraph{A binary pointer tree: few rounds lose almost the whole exponent.}
The pointer-chasing family separates consecutive numbers of rounds, but its
lower bound $2^{\lfloor(b-1)/r\rfloor}/b$ leaves the exponent for a fixed
small number of rounds far below the upper bound $b-r$ of part~(a). The
following contrast closes most of this gap. Let $L\ge1$ and $\beta=L+1$.
Arrange $2^L-1$ \emph{pointer coordinates} as the internal nodes $u$ of a
complete binary tree with levels $0,\ldots,L-1$ and attach $2^L$ \emph{data
coordinates} $Z_v$ to its leaves $v$, so that a block has $2^{\beta}-1$
coordinates. Encode the pointer coordinate of node $u$ as a bit
$x_u\in\{0,1\}$ that selects one of the two children of $u$. The
\emph{pointer path} starts at the root and descends through the selected
children; let $c^\star(x)$ be the leaf it reaches, and
\[
 B_L(x)=Z_{c^\star(x)}=\sum_{v}Z_v\prod_{u\in{\rm path}(v)}\mathbf 1\{x_u\to v\},
\]
where ${\rm path}(v)$ is the set of the $L$ internal nodes above $v$ and
$\{x_u\to v\}$ is the event that the bit of $u$ selects the child on the way
to $v$. Each indicator has degree one, so $\deg B_L=L+1=\beta$, and following
the path with one question per round computes $B_L$ with $\beta$ queries. The
function is the tree evaluation problem of \citet{cook2012}; a deterministic
adversary argument on the same instance shows that exact evaluation with at
most $B$ queries per round needs $L/\lceil\log_2(B+1)\rceil$ rounds
\citep{li2026regex}. Here the total number of queries is capped, exact
evaluation in fewer than $\beta$ rounds is impossible for every strategy, and
we quantify the variance that the missing rounds cost. To reach polynomial
dimension we also take products: for $k\ge1$ let
\[
 B_L^{\otimes k}(x^{(1)},\ldots,x^{(k)})=\prod_{i=1}^kB_L(x^{(i)})
\]
on $k$ disjoint blocks, a contrast of degree $k\beta$ on $k(2^\beta-1)$
coordinates.

\begin{theorem}[Binary pointer trees and their products]\label{thm:binarytree}
Let $L\ge1$, $\beta=L+1$, $k\ge1$ and $b\ge k\beta$, and let
$d=B_L^{\otimes k}$ on $m=k(2^\beta-1)$ coordinates.
\textup{(a)} $V_{\rm branch}(d,b)=1$ and $V_r(d,b)\le2^{b-r}$ for
$1\le r\le b$.
\textup{(b)} For every $1\le r\le b$,
\[
 V_r(d,b)\ \ge\ \frac{2^{k(\beta-1)}}{\big(1+\tfrac{2b}{rk}\big)^{rk}} .
\]
In particular the single tree $B_{b-1}$ on $m=2^b-1$ coordinates ($k=1$,
$\beta=b$) satisfies $V_r(B_{b-1},b)\ge2^{b-1}/(1+2b/r)^r$, and
$V_1(B_{b-1},b)\ge2^{b-1}/b$.
\end{theorem}

\begin{corollary}[Few rounds lose almost the whole exponent]\label{cor:fewrounds}
\textup{(i)} For every $b\ge2$ and $1\le r\le b$,
\[
 \frac{2^{b-1}}{(1+2b/r)^{r}}\ \le\ \sup_{m,d}\frac{V_r(d,b)}{V_{\rm branch}(d,b)}\ \le\ 2^{b-r},
 \qquad\text{in particular}\quad
 \frac{2^{b-1}}{(b+1)^2}\ \le\ \sup_{m,d}\frac{V_2(d,b)}{V_{\rm branch}(d,b)}\ \le\ 2^{b-2}.
\]
Hence for every fixed $r$ the exponent of the $r$-round gap is $b-O(r\log b)$.
\textup{(ii)} An $r$-round protocol that is within a factor $b^{c}$ of full
branching on every contrast of degree at most $b$ needs
$r\log_2(1+2b/r)\ge b-1-c\log_2b$, and therefore $r\ge0.37\,b-(1+c\log_2b)$
rounds: a number of rounds linear in the cap.
\textup{(iii)} Polynomial dimension. For every $b\ge4$ and $r\ge1$, with
$\beta=\lfloor\log_2b\rfloor$ and $k=\lfloor b/\beta\rfloor$, the contrast
$B_{\beta-1}^{\otimes k}$ has degree at most $b$, lives on $m\le b^2$
coordinates, and satisfies $V_{\rm branch}=1$ and
\[
 \log_2V_r\ \ge\ \Big(\frac b\beta-1\Big)(\beta-1)-\frac{rb}{\beta}\log_2\Big(1+\frac{4\beta}{r}\Big)
 \ =\ b-O\Big(\frac{r\,b\log\log b}{\log b}\Big)\quad(b\to\infty,\ r\text{ fixed}).
\]
\end{corollary}
\begin{proof}
(i) The upper bound is Theorem~\ref{thm:roundsfull}(a) and the lower bound is
Theorem~\ref{thm:binarytree}(b) with $k=1$. (ii) Combine $V_r\le b^c$ with
(i) and take logarithms. The function $f(r)=r\log_2(1+2b/r)$ is increasing
and concave, and $f(\alpha^\star b)=b$ for the solution
$\alpha^\star\approx0.376$ of $\alpha\log_2(1+2/\alpha)=1$, where
$f'(\alpha^\star b)=\log_2(1+2/\alpha^\star)-2/((\alpha^\star+2)\ln2)>1.44$. By
concavity $f(r)\le b-1.44(\alpha^\star b-r)$ for $r\le\alpha^\star b$, so
$f(r)\ge b-1-c\log_2b$ forces
$r\ge\alpha^\star b-(1+c\log_2b)/1.44$. (iii) Here $2^\beta\le b$, so
$m=k(2^\beta-1)<kb\le b^2$, and $k\beta\le b<(k+1)\beta\le2k\beta$ gives
$b/k<2\beta$. Theorem~\ref{thm:binarytree}(b) and $k\ge b/\beta-1$,
$k\le b/\beta$ give the display; with $\beta=\lfloor\log_2b\rfloor$ the
subtracted terms are $O(b/\log b)$ and $O(rb\log\log b/\log b)$.
\end{proof}

For two rounds the bound of (i) exceeds one once $2^{b-1}>(b+1)^2$, that is
for $b\ge8$; for smaller caps Appendix~\ref{app:pointerexact} computes the
round constants of $B_{b-1}$ numerically (two rounds lose at least $1.80$ at
$b=3$ and $3.32$ at $b=4$). The single tree uses a number of coordinates
exponential in $b$, like the address contrast; part (iii) removes this at the
price of the lower-order term in the exponent, exactly as products of
address blocks do for fixed forms (Proposition~\ref{prop:productaddress}).
Whether the factor $(1+2b/r)^r$ can be removed, so that the $r$-round gap is
$2^{b-r}$ up to constants, is open.

\paragraph{Proof of Theorem~\ref{thm:binarytree}.}
(a) Following the pointer path of every block returns $d$ itself with
$k\beta\le b$ queries, so $V_{\rm branch}\le1$; every unbiased design has
$\sum_\iota g_\iota(x)=d(x)=\pm1$, and Cauchy--Schwarz gives
$\sum_\iota g_\iota(x)^2/q_\iota\ge1$ at every $x$, so $V_{\rm branch}\ge1$.
The remaining inequality is Theorem~\ref{thm:roundsfull}(a).

(b) Let $\U$ be the uniform law on all coordinates. Under $\U$ the blocks
are independent, $\E_\U[d^2]=1$, and the pointer path of each block is a
uniformly distributed root-to-leaf path. By the dual
formula~\eqref{eq:vdual} with witness $h=d$,
$V_r(d,b)\ge1/\max_{\iota\in\A_r(b)}E_\iota$ with
$E_\iota=\E_\U[\E_\U[d\mid\ell_\iota(X)]^2]$. Fix a deterministic $r$-round
interview $\iota$. Let $S_t$ be the set of coordinates it reads in round $t$;
$S_t$ is a function of the transcript $T_{t-1}$ of rounds $1,\ldots,t-1$, and
we may assume that no coordinate is read twice. For block $i$ write
$a_{t,i}$ and $\lambda_{t,i}$ for the numbers of its pointer and data
coordinates in $S_t$ and $n_{t,i}=a_{t,i}+\lambda_{t,i}$, so that
$\sum_{t,i}n_{t,i}\le b$; let $I_{t,i}$ be the set of pointer coordinates of
block $i$ read in rounds $1,\ldots,t$, ${\rm Obs}_{t,i}$ the set of its leaves
whose data coordinate is read in rounds $1,\ldots,t$, and $c^\star_i$ the leaf
selected in block $i$.

\emph{Step 1: only selected leaves carry signal.}
The event $\{T_r=\tau\}$ fixes the values of the coordinates read under
$\tau$ and nothing else, so conditionally on the transcript the unread
coordinates remain independent and uniform and the blocks remain
independent. Hence
$\E_\U[d\mid\ell_\iota]=\prod_i\E_\U[B_L(X^{(i)})\mid\ell_\iota]$. Within a
block, the data coordinate of an unread leaf $v$ is independent of the
transcript and of the pointer path, so
$\E_\U[Z_v\mathbf 1\{c^\star_i=v\}\mid\ell_\iota]=0$, and
\[
 \big|\E_\U[B_L(X^{(i)})\mid\ell_\iota]\big|
 =\Big|\sum_{v\in{\rm Obs}_{r,i}}Z_v\,\PP(c_i^\star=v\mid T_r)\Big|
 \le\pi_{r,i},\qquad \pi_{t,i}:=\PP(c_i^\star\in{\rm Obs}_{t,i}\mid T_t).
\]
Since $0\le\pi_{r,i}\le1$,
\begin{equation}
 E_\iota\le\E\Big[\prod_{i=1}^k\pi_{r,i}^2\Big]\le\E\Big[\prod_{i=1}^k\pi_{r,i}\Big].\label{eq:allhit}
\end{equation}

\emph{Step 2: the conditional law of a selected leaf.}
For a leaf $v$ of block $i$ let $R_t(v)=|{\rm path}(v)\cap I_{t,i}|$ be the
number of pointer coordinates on its path read by the end of round $t$, and
let $C_t(v)$ be the event that every one of them selects the child towards
$v$. Since $c_i^\star=v$ exactly when all $L$ pointer coordinates on
${\rm path}(v)$ select the child towards $v$, and the unread ones are
uniform given $T_t$,
\begin{equation}
 \PP(c_i^\star=v\mid T_t)=\mathbf 1\{C_t(v)\}\,2^{R_t(v)-L}.\label{eq:leaflaw}
\end{equation}

\emph{Step 3: a potential for one block.}
Put $M_{t,i}=\max_v\mathbf 1\{C_t(v)\}2^{R_t(v)}$ over the leaves of block
$i$, so that $M_{0,i}=1$ and, by \eqref{eq:leaflaw},
$\max_v\PP(c_i^\star=v\mid T_t)=M_{t,i}2^{-L}$, and put
\[
 Q_{t,i}=M_{t,i}+2^L\pi_{t,i},\qquad Q_{0,i}=1 .
\]
We claim
\begin{equation}
 \E[Q_{t,i}\mid T_{t-1}]\le(1+2a_{t,i}+\lambda_{t,i})\,Q_{t-1,i}\le(1+2n_{t,i})\,Q_{t-1,i}.\label{eq:potential}
\end{equation}
\emph{The most likely leaf.} We show
$\E[M_{t,i}\mid T_{t-1}]\le(1+2a_{t,i})M_{t-1,i}$. Condition on $T_{t-1}$,
which fixes $S_t$, drop the block index, and let $v^\star$ attain $M_t$. If no
coordinate of $S_t$ lies on ${\rm path}(v^\star)$, then
$R_t(v^\star)=R_{t-1}(v^\star)$ and $C_t(v^\star)=C_{t-1}(v^\star)$, so
$M_t\le M_{t-1}$. Otherwise let $u$ be the deepest node of
${\rm path}(v^\star)\cap S_t$, write ${\rm path}(u)$ for the nodes above $u$
together with $u$ itself, and define
\[
 G(u)=\mathbf 1\{\text{every node of }({\rm path}(u)\setminus\{u\})\cap I_t\text{ selects towards }u\}\,2^{|{\rm path}(u)\cap I_t|},
\]
\[
 F(u)=\max_{w\text{ below }u}\mathbf 1\{\text{every node of }({\rm path}(w)\setminus{\rm path}(u))\cap I_{t-1}\text{ selects towards }w\}\,2^{|({\rm path}(w)\setminus{\rm path}(u))\cap I_{t-1}|}.
\]
No node of ${\rm path}(v^\star)$ strictly below $u$ is read in round $t$, so
every read node there belongs to $I_{t-1}$ and
$\mathbf 1\{C_t(v^\star)\}2^{R_t(v^\star)}\le G(u)F(u)$: the factor $G(u)$
accounts for the nodes of ${\rm path}(u)$ other than $u$ and for the count
including $u$, the bound $\mathbf 1\{x_u\to v^\star\}\le1$ drops $u$'s own
selection, and $F(u)$ dominates the contribution of the nodes below $u$
because $w=v^\star$ is admissible in its maximum. Consequently
$M_t\le M_{t-1}+\sum_{u\in S_t\cap U}G(u)F(u)$, the sum being over the
pointer coordinates of the block in $S_t$. The quantity $F(u)$ is
$T_{t-1}$-measurable. Given $T_{t-1}$, the $|{\rm path}(u)\cap S_t|-1$ newly
read nodes strictly above $u$ select towards $u$ independently with
probability $1/2$ each, so
\[
 \E[G(u)\mid T_{t-1}]
 =\mathbf 1\{\text{every node of }{\rm path}(u)\cap I_{t-1}\text{ selects towards }u\}\,
   2^{|{\rm path}(u)\cap I_{t-1}|}\cdot2 .
\]
Multiplying by $F(u)$ reassembles, for the leaf $w$ below $u$ attaining
$F(u)$, exactly $2\cdot\mathbf 1\{C_{t-1}(w)\}2^{R_{t-1}(w)}\le2M_{t-1}$,
because $u\notin I_{t-1}$. Summing over the $a_t$ pointer coordinates of the
block in $S_t$ proves the claim.

\emph{The leaves already read.} The sets ${\rm Obs}_{t-1,i}$ and $S_t$ are
$T_{t-1}$-measurable and disjoint, so
$\pi_{t,i}=\PP(c_i^\star\in{\rm Obs}_{t-1,i}\mid T_t)+\PP(c_i^\star\in S_t\mid T_t)$
(we write $c_i^\star\in S_t$ when the data coordinate of that leaf is read in round $t$),
and by the tower property and \eqref{eq:leaflaw},
\[
 \E[\pi_{t,i}\mid T_{t-1}]
 =\pi_{t-1,i}+\PP(c_i^\star\in S_t\mid T_{t-1})
 \le\pi_{t-1,i}+\lambda_{t,i}\,M_{t-1,i}\,2^{-L}.
\]
Adding $2^L$ times this inequality to the bound for $M_{t,i}$ gives
$\E[Q_{t,i}\mid T_{t-1}]\le(1+2a_{t,i}+\lambda_{t,i})M_{t-1,i}+2^L\pi_{t-1,i}$,
which is~\eqref{eq:potential}.

\emph{Step 4: blocks multiply.}
Given $T_{t-1}$ the batch $S_t$ is fixed, the answers of round $t$ in
different blocks are independent, and $Q_{t,i}$ is a function of $T_{t-1}$
and of the round-$t$ answers in block $i$. Hence
$\E[\prod_iQ_{t,i}\mid T_{t-1}]=\prod_i\E[Q_{t,i}\mid T_{t-1}]\le\prod_i(1+2n_{t,i})\prod_iQ_{t-1,i}$.
Let $\Psi_t=\prod_iQ_{t,i}\prod_{s\le t}\prod_i(1+2n_{s,i})^{-1}$. As the
$n_{t,i}$ are $T_{t-1}$-measurable, $\E[\Psi_t\mid T_{t-1}]\le\Psi_{t-1}$ and
$\E[\Psi_r]\le\Psi_0=1$.

\emph{Step 5: combination.}
The pathwise constraint $\sum_{t\le r}\sum_in_{t,i}\le b$ and the inequality
between arithmetic and geometric means over the $rk$ pairs $(t,i)$ give
$\prod_{t\le r}\prod_i(1+2n_{t,i})\le(1+2b/(rk))^{rk}$. Since
$2^L\pi_{r,i}\le Q_{r,i}$,
\[
 \E\Big[\prod_{i=1}^k\pi_{r,i}\Big]\le2^{-Lk}\,\E\Big[\prod_iQ_{r,i}\Big]
 =2^{-Lk}\,\E\Big[\Psi_r\prod_{t\le r}\prod_i(1+2n_{t,i})\Big]
 \le2^{-k(\beta-1)}\Big(1+\frac{2b}{rk}\Big)^{rk}.
\]
With \eqref{eq:allhit} and the dual formula this proves the first statement
of (b). For $r=1$ and $k=1$, \eqref{eq:leaflaw} at $t=0$ gives directly
$\PP(c^\star\in{\rm Obs})\le\lambda_12^{-L}\le b2^{-L}$, which is the last
statement. \hfill$\square$

The argument is a potential-function version of the lucky-hit bound: a
round can raise the conditional probability of the most likely leaf only by
a factor linear in the number of pointer coordinates it reads, whereas
following the path raises it by a factor two per question; the potential
$Q$ adds the probability that a leaf already read is the selected one, which
makes the bound multiplicative across blocks even though an interview may
split its questions among blocks adaptively. As a numerical check, the
released code verifies \eqref{eq:potential} exhaustively on the tree with
$L=2$ (all 4,752 pairs of a state and a nonempty batch; the largest ratio of
the two sides is exactly one) and on random states and batches for $L=3$
and $L=4$ (largest ratio one and $0.83$).


\paragraph{Interpretation.}
Each round can advance the known part of the pointer path by at most one
level unless a question lands on an unknown pointer target, which happens
with probability at most $1/D$ per question. Additional questions within a
round therefore do not substitute for a round. Practical counterparts of
rounds are questionnaire pages, clinic visits and batches of laboratory
orders. The address contrast of Theorem~\ref{thm:universal} is the case
$r=1$, where two rounds already attain the branching coefficient. The
pointer-chasing family shows that this is not a general property of
two-stage designs.

\subsection{Two Rounds Already Attain the Universal Extremum}\label{app:rounds}
An $R$-round interview uses at most $R$ batches, choosing each batch of coordinates using only the preceding batches' answers, with at most $b$ coordinates in total. Let $V_R$ use all such interviews. One round is the fixed-form menu and $b$ rounds suffice for full branching. For any $R\ge2$,
\[
 1\le V_{\rm fixed}/V_R\le K_{m,b},\qquad
 \sup_{m,d\ne0}\frac{V_{\rm fixed}(d,b)}{V_R(d,b)}=2^{b-1},
\]
where the supremum is over identifiable binary contrasts. The upper bound follows from menu inclusion. The address construction attains equality using two batches: all address bits, then the addressed data bit. Thus the universal extremum measures the value of permitting a second batch; it does not require $b$ sequential rounds.

The block construction also needs at most two rounds throughout its identifiable budget curve. At $b=2k+R$, $0\le R\le k$, use the paper's three-query strategy on $R$ blocks and two-query strategy on the others. Ask the first two queries in every block together; then ask the selected third coordinate in the $R$ adaptive blocks. At $b\ge3k$ this already evaluates the contrast exactly. The address and block examples do not separate two rounds from unrestricted branching. Corollary~\ref{cor:roundseparation} supplies a constant separation with a four-feature contrast, and the pointer constructions of Appendix~\ref{app:roundshier} supply the exponential hierarchy.

\subsection{Labels and Across-Person Adaptation}
Let $R_{\rm pre}$ require selecting the interview independently of the current label, and let $R_{\rm label}$ allow observing that label first, while both permit adaptation across people. For an identifiable contrast, Theorem~\ref{thm:asym} and the iid upper construction imply
\[
 \lim_{n\to\infty}nR_{\rm pre}(d,n)
 =\lim_{n\to\infty}nR_{\rm label}(d,n)=V(d,\A).
\]
Thus neither access to earlier people nor early access to the current label reduces the first-order worst-case coefficient relative to an optimal iid design.

There is also an exact statement for the \emph{finite program}. If one extends $\Phi_n$ to separate form laws and score tables for $y=-1,+1$, its objective is the maximum over these labels of their individual worst-profile quadratic objectives. Multiplying the negative-label tables by $-1$ turns both subproblems into the definition of $\Phi_n(d,\A)$. Each has value at least $\Phi_n$, and using the same form law and sign-adjusted tables attains it. Early labels therefore do not improve this program. This is not an equality claim for the exact finite-$n$ minimax risks, which the program bounds only within a constant.

\subsection{Round Constants of the Binary Pointer Tree at Small Caps}\label{app:pointerexact}
The two-round bound of Theorem~\ref{thm:rounds}(c) exceeds one only for $b\ge8$. For small caps we compute the round constants of $B_{b-1}$ directly. Under the uniform law $\U$ with a label independent of the profile, $M_r(\U)$, the optimal second moment over $r$-round designs, is a lower bound on $V_r$, because $V_r=\sup_PM_r(P)$ by Sion's theorem applied to \eqref{eq:unbiasedprogram}. For $b\le4$ the cube is small enough to evaluate the uniform-optimal design at every profile; its worst profile is an upper bound on $V_r$. The solver of Appendix~\ref{app:engine} computes $M_r(\U)$ with certified bounds that agree to $10^{-7}$ or better.
\begin{table}[h]
\centering
\caption{Round constants of the binary pointer tree $B_{b-1}$ ($m=2^b-1$). Entries are $M_r(\U)$, a lower bound on $V_r$; the brackets give the worst profile of the uniform-optimal design, an upper bound on $V_r$, when it differs. The last columns are the upper bound $2^{b-2}$ of Theorem~\ref{thm:rounds}(a) and the lower bound $2^{b-1}/(b+1)^2$ of Theorem~\ref{thm:rounds}(c) for $r=2$.}
\label{tab:pointer}
\small
\begin{tabular}{rrllllrr}
\toprule
$b$ & $m$ & fixed & $r=2$ & $r=3$ & $r=4$ & $2^{b-2}$ & $2^{b-1}/(b+1)^2$ \\
\midrule
3 & 7 & 4.000 & 1.800 & 1.000 &  & 2 & 0.25 \\
4 & 15 & 8.000 & 3.324 [3.334] & 1.800 & 1.000 & 4 & 0.32 \\
\bottomrule
\end{tabular}
\end{table}

At $b=3$ and $b=4$ the two bounds agree to the reported precision except for two rounds at $b=4$, where they differ by 0.3\%. Two rounds therefore lose almost the factor $2^{b-2}$ allowed by Theorem~\ref{thm:rounds}(a) already at these caps, and each missing round costs a factor close to two. The cap $b=5$ ($m=31$, 169,911 forms) was not completed.


